% Part: computability
% Chapter: computability-theory
% Section: halting-problem

\documentclass[../../../include/open-logic-section]{subfiles}

\begin{document}

\olfileid{cmp}{thy}{hlt}
\olsection{د درېدنې مسئله}

د جوړښت له مخې نړيواله جزوي محاسبه کېدونکې تابع~$\fn{Un}(e,x)$ هله
او يوازې هله تعريف شوې ده چې په~$e$ کوډ شوې د تابع محاسبه د ننوت~$x$
دپاره يو قيمت پيدا کړي۔ طبيعي پوښتنه دا ده چې ايا موږ پرېکړه کولے شو
چې داسې ده که نه۔ په حقيقت کښې دا پرېکړه نۀ شو کولے۔ د محاسبې د
ټيورينګ ماشين په مدل کښې دا په دې مانا ده چې ايا ورکړل شوے ټيورينګ
ماشين پر ورکړل شوي ننوت درېږي که نه، په محاسبوي ډول د پرېکړې وړ نۀ
ده۔ له همدې کبله لاندې قضيه د ``درېدنې د مسئلې د ناپرېکړتيا'' په نوم
پېژندل کېږي۔ لاندې به دوه ثبوتونه ورکړو۔ لومړے زموږ د مخکنۍ بحث لړۍ
پسې غځوي، او دويم ډېر نېغ دے۔

\begin{thm}
\ollabel{thm:halting-problem}
فرض کړئ
\[
h(e, x)  =
\begin{cases}
1 & \text{که $\fn{Un}(e, x)$ تعريف شوې وي} \\
0 & \text{که داسې نۀ وي۔}
\end{cases}
\]
نو $h$ محاسبه کېدونکې نۀ ده۔
\end{thm}

\begin{proof}
فرض کړئ $h$ محاسبه کېدونکې ده۔ ښيو چې په دې سره يوه نړيواله محاسبه
کېدونکې تابع تعريفولے شو۔ داسې تعريف کړئ:
\[
\fn{Un'}(e,x) =
\begin{cases}
\fn{Un}(e,x) & \text{که $h(e,x) = 1$ وي} \\
0 & \text{که داسې نۀ وي۔}
\end{cases}
\]
خو اوس $\fn{Un'}(e, x)$ هرځاے تعريف شوې تابع ده، او که~$h$ محاسبه
کېدونکې وي نو دا هم محاسبه کېدونکې ده۔ د بېلګې په توګه، د بنسټيز
بازګښت په کارولو سره~$g$ داسې تعريفولے شو:
\begin{align*}
g(0, e, x) & \simeq 0 \\
g(y+1, e, x) & \simeq \fn{Un}(e,x);
\end{align*}
بيا
\[
\fn{Un'}(e,x) \simeq g(h(e,x),e,x).
\]
څرنګه چې $\fn{Un'}(e,x)$ په هر هغه ځاے کښې له~$\fn{Un}(e,x)$ سره
برابره ده چې وروستۍ تابع تعريف شوې وي، نو~$\fn{Un'}$ د هغو جزوي
محاسبه کېدونکو تابعو دپاره نړيواله ده چې هرځاے تعريف شوې وي۔ خو دا له
\olref[nou]{thm:no-univ} سره تناقض دے۔
\end{proof}

\begin{proof}
فرض کړئ $h(e,x)$ محاسبه کېدونکې ده۔ تابع~$g$ داسې تعريف کړئ:
\[
g(x) =
\begin{cases}
  0                & \text{که $h(x,x) = 0$ وي} \\
  \fundefined & \text{که داسې نۀ وي۔}
\end{cases}
\]
تابع~$g$ جزوي محاسبه کېدونکې ده۔ د بېلګې په توګه، دا د
$\umin{y}{h(x,x) = 0}$ په توګه تعريفولے شو۔ نو د کوم~$e$ دپاره د هر
~$x$ په نسبت $g(x) \simeq \fn{Un}(e,x)$ ده۔ ايا~$g$ پر~$e$ تعريف
شوې ده؟ که تعريف شوې وي، نو د~$g$ د تعريف له مخې $h(e,e) = 0$ دے
(کله چې~$g$ تعريف شوې وي، يوازې د~$0$ قيمت اخلي)۔ د~$h$ د تعريف له
مخې دا په دې مانا ده چې $\fn{Un}(e, e)$ تعريف شوې نۀ ده۔ زموږ د دې
فرض له مخې چې د هر~$x$ دپاره $g(x) \simeq \fn{Un}(e, x)$ ده،
$g(e)$ هم تعريف شوې نۀ ده، چې تناقض دے۔ بل پلو، که $g(e)$ تعريف شوې
نۀ وي، نو $h(e,e) \neq 0$، او ځکه $h(e,e) = 1$ دے۔ له دې څخه راځي
چې $\fn{Un}(e, e)$ تعريف شوې ده۔ خو څرنګه چې
$g(x) \simeq \fn{Un}(e, x)$ ده، نو $g(e)$ به هم تعريف شوې وه۔ بيا
تناقض راغے۔
\textbf{د ژباړې د نسخې سمون (OLCMP-022):} سرچينې په لومړي حالت کښې
د قوس دننه دا ځانګړنه د درېدنې تابع ته منسوبه کړې وه، حال دا چې هغه
د قضيې له تعريفه هرځاے تعريف شوې ده۔ د استدلال اړوند خبره دا ده چې
جوړه شوې قطري تابع، کله تعريف شوې وي، يوازې صفر راګرځوي۔ دلته همدا
فاعل سم شوے دے؛ د ثبوت فورمولونه نۀ دي بدل شوي۔
\end{proof}

\begin{tagblock}{TMs}
\begin{explain}
دا استدلال د ټيورينګ ماشينونو په ژبه هم بيانولے شو۔ فرض کړئ داسې
ټيورينګ ماشين~$H$ شته چې د يو ټيورينګ ماشين~$E$ تشريح او يو
ننوت~$x$ اخلي، او پرېکړه کوي چې~$E$ پر~$x$ درېږي که نه۔ بيا به يو
بل ټيورينګ ماشين~$G$ جوړولے شو چې يو ننوت~$x$ اخلي، $H$ چلوي څو
پرېکړه وکړي چې د شاخص~$x$ ماشين~$M_x$ پر ننوت~$x$ درېږي که نه، او
بيا د هغې مخالف کار کوي۔ په بله وينا، که~$H$ ووايي چې~$M_x$ پر~$x$
درېږي، $G$ په ناپايه کړۍ کښې ننوځي؛ او که~$H$ ووايي چې~$M_x$ پر~$x$
نۀ درېږي، نو~$G$ بس درېږي۔ ايا~$G$ د خپل شاخص په ننوت درېږي؟ پورته
استدلال ښيي چې دا هله او يوازې هله درېږي چې نۀ درېږي—يو تناقض۔ نو
زموږ دا فرض چې داسې ټيورينګ ماشين~$H$ شته، بايد ناسم وي۔
\end{explain}
\end{tagblock}

\end{document}
